\newpage
\section{Pseudocode and Schematic Workflow}
\setlength{\textfloatsep}{0pt}
\begin{algorithm}[h]
\small
\SetAlgoLined
\DontPrintSemicolon

\KwIn{Query $q$; early-fused graph $G_d$; edge retriever $f_e$; edge reranker $g_e$;
      node reranker $g_n$; expanded-query retrievers
      $f_{P\!\rightarrow S},\,f_{S\!\rightarrow P}$;
      star-based LLM Refiner $\mathcal{R}_\star$; beam width $b$; final list size $k$.}
\KwOut{Ranked edge list $E_q$.}

$\mathcal{E}_{\text{cand}} \leftarrow f_e(q,G_d)$ \tcp*[l]{retrieve top-$k_1$ edges}
$\mathcal{E}_{\text{cand}} \leftarrow g_e(q,\mathcal{E}_{\text{cand}})$ \tcp*[l]{rerank $\!\rightarrow$ top-$k_2$}
$G_c \leftarrow (V_c,\mathcal{E}_{\text{cand}})$ \tcp*[l]{candidate subgraph}

\BlankLine
$S \leftarrow \text{top-}b$ nodes by $g_n(q,V_c)$ \tcp*[l]{seed selection}
\ForEach{$u\in S$}{
  $q_u \leftarrow [\,q;\Gamma(u)\,]$\;
  \uIf{$u$ is a passage}{
    $N_u \leftarrow f_{P\!\rightarrow S}(q_u,G_d)$
  }\Else{
    $N_u \leftarrow f_{S\!\rightarrow P}(q_u,G_d)$
  }
  Add edges $\{(u,v)\mid v\in N_u\}$ to $G_c$\;
}

$G_\ell \leftarrow G_c$\;
isAgg $\leftarrow \mathcal{R}_\star.\textsc{DetectAgg}(q)$\;

Decompose $G_\ell$ into stars $\{\mathcal{S}_i\}$\;

$G_q \leftarrow \varnothing$\;
\ForEach{$\mathcal{S}_i$}{
  $G_q \leftarrow G_q \,\cup\,
          \mathcal{R}_\star.\textsc{Refine}\!\bigl(\mathcal{S}_i,\mathcal{G}_d,\,
          \texttt{agg}=\,\text{isAgg}\bigr)$\;
}
$E_q \leftarrow$ top-$k$ edges in $G_q$ (sorted by score)\;
\Return $E_q$
\caption{\textsc{HELIOS} Inference Pipeline}
\label{alg:helios_inference}
\end{algorithm}

\begin{figure*}[t]
\centering
    \includegraphics[width=\textwidth]{Figures/helios_workflow.pdf}
    \caption{
        High-level schematic workflow of \texttt{HELIOS}.
    }
\label{fig:helios_workflow}
\end{figure*}
{
\feditcolor
To enhance the clarity of our framework, we provide both a pseudocode and a schematic workflow. 
The pseudocode in Algorithm~\ref{alg:helios_inference} details the step-by-step inference process, encompassing edge-based bipartite subgraph retrieval, query-relevant node expansion, and star-based LLM refinement. 
Additionally, Figure~\ref{fig:helios_workflow} offers a high-level schematic view of the overall workflow.
}
% \setlength{\textfloatsep}{0pt}
% \begin{algorithm}[ht]
% \small
% \SetAlgoLined
% \DontPrintSemicolon

% \KwIn{Training triples $(q,a)$ where $a$ contains the answer edge(s).}
% \KwOut{Fine-tuned modules $f_e,\,g_e,\,g_n,\,f_{P\!\rightarrow S},\,f_{S\!\rightarrow P},\,\mathcal{R}_\star$.}

% \textbf{Edge-level training}\;
% \Indp
% Construct positive \& hard-negative edge sets\;
% Train $f_e$ with contrastive + KD loss (teacher $g_e$)\;
% Fine-tune $g_e$ on the same triples\;
% \Indm

% \BlankLine
% \textbf{Node-level training}\;
% \Indp
% Derive positive/negative node sets from edge triples\;
% Fine-tune node reranker $g_n$\;
% \Indm

% \BlankLine
% \textbf{Expanded-query retrievers}\;
% \Indp
% Build triples $\bigl([q;\Gamma(u)],\,v^+,\,v^-\bigr)$\;
% Train $f_{P\!\rightarrow S}$ and $f_{S\!\rightarrow P}$ with ColBERT loss\;
% \Indm

% \BlankLine
% \textbf{LLM refinement}\;
% \Indp
% Synthesize instruction data (agg.\ detection, column aggregation, passage verification)\;
% Fine-tune $\mathcal{R}_\star$ on the synthetic set\;
% \Indm

% \caption{\textsc{HELIOS} Training Procedure}
% \label{alg:helios-training}
% \end{algorithm}

\section{Training Scheme}
\label{sec:training_scheme}

\subsection{Edge Retriever and Reranker}
\label{sec:edge_retriever_and_reranker}


The training scheme for our encoder $f_{e}$ follows the methodology outlined in \texttt{ColBERTv2} (\citealt{ColBERTv2}), leveraging a combination of in-batch negative loss and knowledge distillation loss to train the model. 
Specifically, the in-batch negative loss treats the edges corresponding to other queries within the same batch as negative samples.
This approach calculates a contrastive loss between the positive and negative edges.
In constructing the training dataset, it is crucial to have both positive and negative edges for each query. 
To define the positive edge, we use passages containing the answer and the associated table segments as ground truth and denoted as $x_{gt}$.
Conversely, negative edges are constructed by combining hard negative tables and passages from prior work (\citealt{COS}) with in-batch negative edges and are denoted as $n(q)$.
The contrastive loss $L_{cl}$ is represented as follows:
\begin{equation}
\small
    L_{cl} = - \sum_{(q,x_{gt})}\log{\frac{exp(s(q,x_{gt}))}{exp(s(q,x_{gt}))+\sum_{z\in n(q)}exp(s(q,z))}}
\end{equation}
\label{eqn:contrastive_loss}
The knowledge distillation process refines the edge encoder using a teacher-student model setup.
The distillation loss is computed based on the KL divergence between the score distribution generated by the teacher model and the training encoder. The training was conducted for 1 epoch with a batch size of 512 and a learning rate of 1e-5. We employed a cosine learning rate scheduler with 40 warm-up steps.

Here, the teacher model is the all-to-all interaction reranker $g_{e}$ fine-tuned with the contrastive loss, which serves as a more precise reference for edge relevance.
This method ensures that the encoder learns from a more sophisticated model, improving its capacity to accurately rank edges based on the query. The training was conducted for 1 epoch with a batch size of 96 and a learning rate of 2e-4. We employed a cosine learning rate scheduler with a warmup ratio of 0.1.



\subsection{Node Reranker}
\label{sec:node_reranker}

The training method for the node reranker $g_n$ is identical to that of the edge reranker $g_e$. 
For constructing the training dataset, we utilize the \texttt{OTT-QA} dataset \citep{OTT-QA}. 
Positive nodes are defined as those directly connected to the nodes that contain the correct answer in \texttt{OTT-QA}. 
In contrast, negative nodes are selected from the set of nodes retrieved through edge-based bipartite subgraph retrieval, excluding any nodes connected to the answer-containing nodes. The training was conducted for 1 epoch with a batch size of 96 and a learning rate of 2e-4. We used a cosine learning rate scheduler with a warm-up ratio of 0.1.


\subsection{Expanded Query Retrievers}
\label{sec:expanded_query_retrievers}

The training scheme for our expanded query retrievers $f_{S\rightarrow P}$, $f_{P\rightarrow S}$ also follow the methodology outlined in \texttt{ColBERTv2} (\citealt{ColBERTv2}).
% To enhance the accuracy of subgraph augmentation, we fine-tuned two separate retrievers: 
% (i) $f_{S\rightarrow P}$ that retrieves passage nodes based on queries expanded with table segment nodes, and
% (ii) $f_{P\rightarrow S}$another that retrieves table segment nodes using passage-expanded queries.
To construct the training dataset, we generated triples consisting of the expanded query, positive node, and negative node. 
Expanded queries were created by incorporating nodes that are connected to the node containing the answer. 
Positive nodes consist of the nodes that contain the answer. Negative nodes are constructed using hard negative nodes as outlined in prior work (\citealt{COS}). The training was conducted for 1 epoch with a batch size of 512 and a learning rate of 1e-5. We employed a cosine learning rate scheduler with a 40 warmup steps.

% \subsection{Expanded Query Retrievers}
% \label{sec:expanded_query_retrievers}

\subsection{Star-based LLM Refinement}
\label{sec:star_based_llm_refinement_finetune}
We fine-tuned the Llama-3.1-Instruct 8B model \cite{dubey2024llama} using an instruction dataset synthesized with GPT-4o \cite{hurst2024gpt}. The dataset includes tasks for aggregate query detection, column-wise aggregation, and passage verification, and training was conducted for 1 epoch with a batch size of 128 and a learning rate of 2e-5. We employed a linear learning rate scheduler with a warmup ratio of 0.03.

To construct the training data, we applied edge-based bipartite graph retrieval and query-relevant node expansion on the OTT-QA training set, providing the results to GPT-4o for output generation. The final dataset was formed by combining these outputs with zero-shot prompts. Due to API cost constraints, we sampled 25\% of the OTT-QA training dataset, resulting in 1,655 samples each for query detection and column-wise aggregation and 7,010 samples for passage verification.
% 우리는 llama-3.1-instruct 8B \cite{dubey2024llama} model을 instruction finetune했다 using a training dataset synthesized with GPT-4o \cite{hurst2024gpt}, incorporating aggregate query detection, column-wise aggregation, and passage verification tasks for 1 epochs.
% OTT-QA의 training data에 대해 edge-based bipartite graph retrieval과 query-relevant node expansion을 수행한 결과를 GPT-4o에게 제공하여 출력을 생성하도록 했고, 생성된 출력과 zero-shot prompt를 합쳐서 training dataset을 구축했다. GPT-4o의 API Cost로 인해 전체 OTT-QA training dataset의 25%만 sample한 training dataset을 구축했고, 그 결과 query detection 1,655, column-wise aggregation 1,655, passage verification 7,010개로 training dataset을 구성했다.
% The batch size is set to 128 and the learning rate is set to 2e-5. We use the linear learning rate scheduler policy with the warmup ratio set to 0.03. 
% The training is conducted on 4 A6000 GPUs for 4 hours.



\section{Star-based LLM Refinement Supplementary}
\label{sec:star_based_llm_refinement_supp}
\begin{figure*}
\centering
    \includegraphics[width=\textwidth]{Figures/Star-based_LLM_Refinement.pdf}
    \caption{
        The overall process of star-based LLM refinement for queries classified as aggregation queries. 
        Table segment nodes of the same color (black, orange) indicate segments that belong to the same original table.
    }
\label{fig:star_based_llm_refinement}
\end{figure*}

In the Star-based LLM Refinement stage, the expanded graph from the Query-relevant Node Expansion (QNE) step undergoes further enhancement through the reasoning capabilities of LLMs. This stage consists of two primary operations: (1) aggregation of query-relevant table segments and (2) verification of query-irrelevant passages. The detailed process is illustrated in Figure~\ref{fig:star_based_llm_refinement}.
















\section{Experiment Supplementaries}
\label{sec:experiment_supp}

\subsection{Implementation Details}
In our edge generation step (\textsection~\ref{sec:candidate_bipartite_graph_generation}), we used the same named entity recognition and entity linking models used by \texttt{COS} (\citealp{COS}). 
For the late-interaction edge retriever $f_e$ (\textsection~\ref{sec:candidate_bipartite_graph_generation}) and the expanded query retrievers $f_{P \rightarrow S}$ and $f_{S \rightarrow P}$ (\textsection~\ref{sec:candidate_bipartite_graph_augmentation}), we employed \texttt{ColBERTv2} (\citealp{ColBERTv2}) as the baseline model.
For the all-to-all interaction edge reranker $g_e$ (\textsection~\ref{sec:candidate_bipartite_graph_generation}) and node reranker $g_n$ (\textsection~\ref{sec:candidate_bipartite_graph_augmentation}), we used the \texttt{bge-reranker-v2-minicpm-layerwise} (\citealp{bge-reranker-v2-minicpm-layerwise}), specifically utilizing layer 24 as the baseline model. 
Lastly, for star-based LLM refinement (\textsection~\ref{sec:graph_refinement_via_complex_reasoning}), we used \texttt{Llama-3.1-8B-Instruct} (\citealp{dubey2024llama}) as the large language model.
In our experiments, the value of $k_1$ for the edge retriever $f_e$ was set to 400. 
Since \texttt{COS} selects the top-200 nodes as seed nodes, we fixed $k_2$ for the edge reranker $g_e$ to 100 to ensure a fair comparison.
We evaluated \texttt{HELIOS} on the OTT-QA's development set (2,214 examples) using four A100-80GB GPUs, which required approximately 3.4 hours.

\subsection{Parameter Sensitivity of Retriever}
\label{sec:exp_parameter_sensitivity}

\begin{figure}
\centering
    \includegraphics[width = \columnwidth]{Figures/beam_size_sensitivity.pdf}
    \caption{Change in AR@50 with varying beam width}
\label{fig:beam_width}
\end{figure}

We explored the impact of varying beam width $b$ on retrieval accuracy in terms of AR@50. 
The beam width directly influences the number of expanded nodes (\textsection~\ref{sec:candidate_bipartite_graph_augmentation}).
We experimented with beam widths of 0, 2, 5, 10, 25, 50 and measured the corresponding changes in AR@50.

Figure~\ref{fig:beam_width} illustrates the change in AR@50.
We observed that AR@50 was improved by 1.7\% as the beam width monotone increased from 0 to 10.
This indicates that larger beam widths lead to more accurate node augmentations by performing a more exhaustive search across the expanding node space. 
Interestingly, when the beam size increased to 50, AR@50 decreased slightly by 0.4\% compared to beam size 10. 
This drop may be due to LLM hallucinations in the star-based LLM refinement (SLR) module, where irrelevant edges were added to $G_l$, causing the SLR to fail in selecting the correct query-relevant nodes.
This observation highlights the importance of selectively expanding only the most probable nodes within the query-relevant node expansion module.

\subsection{Parameter Sensitivity of Reader}
\label{sec:reading_accuracy_supp}
\begin{figure}[t]
    \centering
    \small
    \begin{tabular}{c@{}c@{}}
        \includegraphics[width=0.5\columnwidth]{Figures/Plots/GPT4o/dev_em.pdf} & 
        \includegraphics[width=0.5\columnwidth]{Figures/Plots/GPT4o/dev_f1.pdf} \\
        \noalign{\vskip -1mm}
        % (a) EM Score (GPT-4o) & (b) F1 Score (GPT-4o) \\
        \multicolumn{2}{c}{(a) GPT-4o} \\
        \includegraphics[width=0.5\columnwidth]{Figures/Plots/Llama31-Instruct-70B/dev_em.pdf} & 
        \includegraphics[width=0.5\columnwidth]{Figures/Plots/Llama31-Instruct-70B/dev_f1.pdf} \\
        \noalign{\vskip -1mm}
        % (c) EM Score (Llama3.1-Instruct-80B) & (d) F1 Score (GPT-4o) \\
        \multicolumn{2}{c}{(b) Llama3.1-Instruct-70B} \\
        % \includegraphics[width=0.5\columnwidth]{Figures/Plots/GPT4o/dev_em.pdf} & 
        % \includegraphics[width=0.5\columnwidth]{Figures/Plots/GPT4o/dev_f1.pdf} \\
        % \noalign{\vskip -1mm}
        % % (e) EM Score (Reader 3) & (f) F1 Score (Reader 3) \\
        % \multicolumn{2}{c}{(c) FiE-330M (Fine-tuned)} 
        % \\    
    \end{tabular}
    \vspace{-3mm}
    \caption{
        End-to-end QA accuracy comparison across different top-$k$ values on the OTT-QA dev set.
    }
    \label{fig:em_f1_comparison_over_diff_k}
\end{figure}
% Figure~\ref{fig:em_f1_comparison_over_diff_k} presents the end-to-end QA accuracy of \texttt{HELIOS}, \texttt{COS}, and \texttt{COS w/ ColBERT \& bge} across different reader models, including \texttt{Llama-3.1-70B-Instruct} (\citealp{dubey2024llama}) and \texttt{GPT-4o} (\citealp{hurst2024gpt}), and different $k$ values, where $k\in\{2, 5, 10, 20, 50\}$.
Figure~\ref{fig:em_f1_comparison_over_diff_k} presents the end-to-end QA accuracy of \texttt{HELIOS}, \texttt{COS}, and \texttt{COS w/ ColBERT \& bge} across different reader models—\texttt{Llama-3.1-70B-Instruct} (\citealp{dubey2024llama}) and \texttt{GPT-4o} (\citealp{hurst2024gpt})—with varying top-$k$ values ($k \in {2, 5, 10, 20, 50}$). 
Due to budget constraints, we sampled 10\% (221 out of 2,214) of the development set's QA pairs when evaluating accuracy variations for \texttt{GPT-4o}. 
As a result, \texttt{HELIOS} consistently achieves the highest AR@$k$ across all $k$ values. 
It improves the average EM score by 15.5\% for \texttt{Llama-3.1-70B-Instruct} and 9.2\% for \texttt{GPT-4o} compared to \texttt{COS w/ ColBERT \& bge}.

\subsection{Error and Reader Impact Analysis}
\label{sec:reader_impact}
{
\feditcolor
While \texttt{HELIOS} achieves high retrieval accuracy, there remains room for improvement in end-to-end QA performance. 
To investigate the gap between retrieval and QA accuracy, we performed an error analysis on 408 OTT-QA questions for which \texttt{HELIOS} successfully retrieved the gold evidence but still produced incorrect answers (F1 = 0).
We then re-tested these questions using \texttt{DeepSeek R1}~\citep{liu2024deepseek}, a stronger long-context reasoning model, with the same retrieved evidence.
Remarkably, the average F1 improved from 0 to 29.8, confirming that a more capable reader can better leverage the available evidence, especially for multi-hop questions, and further reduce hallucinations.
These findings highlight the potential for additional gains if a more advanced QA model is used on top of \texttt{HELIOS}’s strong retrieval.
}


\subsection{Detailed Breakdown of Offline Overhead}
\label{sec:appendix_offline_overhead}

{\feditcolor
Our method involves an offline stage for graph construction and retriever fine-tuning. 
The entire offline pipeline took approximately 49 hours, with graph construction accounting for about 38 hours. 
We did not optimize this phase since it is not the focus of our work. 
However, this stage can be significantly accelerated via straightforward parallelization (e.g., multi-core or multi-node setups), meaning the overhead does not compromise the feasibility of our approach.

Table~\ref{tab:graph_construction_time} provides a detailed breakdown of the graph construction process, which alone accounted for about 38 hours. 
In this phase, we construct a cross-modal retrieval graph by linking 400K tables and 5M passages via entity recognition and linking, following the approach of COS~\citep{COS}.

\begin{table}[h]
    \centering
    \setlength{\tabcolsep}{2pt}
    \resizebox{0.8\columnwidth}{!}{%
        \begin{tabular}{l|r}
        \hline
        \textbf{Component} & \textbf{Processing Time} \\
        \hline
        Entity recognition & 3h 14m 32s \\
        Entity linking & 2h 20m 48s \\
        Edge indexing & 22h 29m 54s \\
        Passage indexing & 7h 49m 01s \\
        Table segment indexing & 2h 29m 09s \\
        \hline
        \end{tabular}
    }
    \vspace{-3mm}
    \caption{Graph construction time by component.}
    \label{tab:graph_construction_time}
\end{table}

Table~\ref{tab:retriever_finetuning_time} summarizes the time required to fine-tune each retriever and reranker submodule used in our system. 
This phase took around 10 hours in total and can be parallelized with multi-core or multi-node setups. 
Given that both graph construction and retriever fine-tuning are offline processes, the overhead does not hinder the practicality of our system.

\begin{table}[h]
    \centering
    \setlength{\tabcolsep}{2pt}
    \resizebox{\columnwidth}{!}{%
        \begin{tabular}{l|r}
        \hline
        \textbf{Component} & \textbf{Training Time} \\
        \hline
        Edge retriever & 2h 49m 37s \\
        Edge reranker & 1h 25m 32s \\
        Node reranker & 1h 39m 20s \\
        Expanded query retriever (passages) & 36m 32s \\
        Expanded query retriever (tables) & 4h 14m 45s \\
        \hline
        \end{tabular}
    }
    \vspace{-3mm}
    \caption{Fine-tuning time for retriever submodules.}
    \label{tab:retriever_finetuning_time}
\end{table}
}


% HELIOS consistently outperforms all competitors on AR@$k$ across different $k$ values.
% \texttt{HELIOS} improved the performance of all $k$, achieving an average F1 score improvement of 9.0\% and 8.7\% compared to \texttt{COS w/ ColBERT \& bge} for \texttt{Llama-3.1-70B-Instruct} and \texttt{GPT-4o} each.  
% Figure~\ref{fig:em_f1_comparison_over_diff_k} compares the end-to-end QA accuracy of \texttt{HELIOS}, \texttt{COS}, and \texttt{COS w/ ColBERT \& bge} across different reader models—\texttt{Llama-3.1-70B-Instruct} (\citealp{dubey2024llama}) and \texttt{GPT-4o} (\citealp{hurst2024gpt})—and varying top-$k$ values, where $k \in {2, 5, 10, 20, 50}$.
% \texttt{GPT-4o}의 경우 We did not have enough budget to conduct multiple runs for our experiments using GPT-4o API이므로 $k$에 따른 정확도 변화를 측정할 때는 development set의 2214개의 QA pairs 중 10%인 221개만 random sample하여 사용했다.
% \texttt{HELIOS} consistently achieves the highest AR@$k$ across all values of $k$. 
% Specifically, it improves performance at every retrieval depth, yielding an average F1 score increase of 9.0\% and 8.7\% over \texttt{COS w/ ColBERT \& bge} for \texttt{Llama-3.1-70B-Instruct} and \texttt{GPT-4o}, respectively.
% \subsection{Trade-off Between Computational Overhead and Performance Gains}
% {\editcolor 
% \subsection{Algorithm Execution Time}
% \label{sec:exp_algorithm_execution_time}

% \begin{table}[t]
%     \centering
%     % \setlength{\tabcolsep}{2pt} % 열 간격 조정
%     \resizebox{\columnwidth}{!}{%
%     \begin{tabular}{l|c|c}
%         \hline
%         \textbf{Algorithm}      & \textbf{Execution Time (s)}     & \textbf{nDCG@50}  \\
%         \hline
%         \texttt{CORE}  & 4.13     & 25.4 \\
%         \texttt{COS}  & 3.75     & 33.6 \\
%         \texttt{COS w ColBERT \& bge}  & 5.46     & 36.5 \\
%         \texttt{HELIOS}  & 5.14     & 47.0 \\
%         \texttt{HELIOS w Finetuned SLR} & 4.76 & \textbf{47.6} \\
%         \hline
%         \texttt{w/o SLR}                 & \textbf{2.16}             & 46.5 \\
%         \hline
%     \end{tabular}
%     }
%     \caption{\label{tab:computational_overhead}
%         Algorithm execution time comparison of HELIOS against competitors
%     }
% \end{table}

% Table~\ref{tab:computational_overhead} compares the execution time of \texttt{HELIOS} against \texttt{COS} and \texttt{CORE}. 
% All experiments were conducted under identical conditions using four A100-80GB GPUs, with LLM inference handled by the SGLang inference engine.

% The results demonstrate that \texttt{HELIOS} effectively balances retrieval speed and accuracy. 
% Compared to \texttt{COS}, it increases execution time by 37.1\% (24.5\% higher than \texttt{CORE}) while improving nDCG@50 by 39.9\% (85.0\% higher than \texttt{CORE}). 
% Additionally, leveraging a fine-tuned SLR reduces execution overhead to 26.9\% while further enhancing nDCG@50 by 41.7\% compared to \texttt{COS}. 
% This improvement stems from a more efficient design that eliminates the need for few-shot examples.

% Without SLR, \texttt{HELIOS} runs 42.4\% faster than \texttt{COS} while achieving a 38.4\% improvement in nDCG@50. 
% Despite incorporating ColBERT, \texttt{HELIOS} maintains a speed advantage over \texttt{COS}, as \texttt{COS} performs expanded query retrieval on all 200 retrieved table segments, whereas \texttt{HELIOS} employs a beam search algorithm, restricting retrieval to a beam width of 10.
% Furthermore, modifying \texttt{COS} to incorporate ColBERT retrieval and the bge reranker increased its execution time by 45.6\% compared to its original implementation.
% }

\begin{figure*}
\centering
    \includegraphics[width = \textwidth]{Figures/Qualitative_analysis.pdf}
    \caption{
        Qualitative analysis on four question-answer pairs.
        (a) A case where passage verification is successful.
        (b) A first case where passage verification has failed.
        (c) A second case where passage verification has failed.
        (d) A case where table aggregation is successful.
    }
\label{fig:qualitative_analysis}
\end{figure*}

\section{Qualitative Analysis}
\label{sec:qualitative_analysis}


In this section, we present a qualitative analysis of HELIOS's Column-wise Aggregation module and Passage Verification module, with the results illustrated in Figure~\ref{fig:qualitative_analysis}.
The subfigures in Figure~\ref{fig:qualitative_analysis} showcase the performance and distinctive scenarios for each module: 
(a) highlights successful cases of the Column-wise Aggregation module, while (b), (c), and (d) demonstrate representative cases related to the Passage Verification module. 
For each subfigure, the query is depicted in dark blue, the data provided to the sub-module are shown in light blue, and the inference result from the LLM is encapsulated in a purple speech bubble with a llama icon.


Figure~\ref{fig:qualitative_analysis}(a) shows a successful case of the column-wise aggregation module in resolving a complex query: identifying the birth date of the "most recent Segunda Liga Player of the Month." 
The essential part of answering this question was to recognize that the most recent player, Basilio Almeida, was honored in November 2009, as indicated in the SJPF Segunda Liga Player of the Month table. 
However, the initial data lacked the table segment containing the relevant row. 
The column-wise aggregation module reconstructed the table as shown in Figure~\ref{fig:qualitative_analysis}(a) to include this missing information, enabling the system to restore the row with the necessary details. 
The LLM correctly inferred from the reconstructed table that the row corresponding to the most recent player was Row 4, based on the Year and Month columns. 
This led \texttt{HELIOS} to accurately generate the final answer in this question, which is "12 August 1971."


Figure~\ref{fig:qualitative_analysis}(b) shows a successful case of the passage verification module in addressing the query, "How many years did the series that Zuzanna Szadkowski appeared in for 3 episodes run for?".
The module was provided with a Zuzanna Szadkowski table summarizing her appearances and a set of associated passages. 
The "Notes" column of the table segment confirmed that she appeared in three episodes of the series Guiding Light. 
The module correctly identified the one mentioning Guiding Light among the provided passages, the one which indicated that the series was broadcast on CBS for 57 years. 
the module correctly verified that the passage using the passage's information noting its broadcast duration, leading to an accurate answer.


Figure~\ref{fig:qualitative_analysis}(c) shows a failure case of the Passage Verification module when answering the query, "What is the province where the unit in the Morgan District Brigade that disbanded in 1782 was founded?".
The module correctly identified 'Burke County Regiment' as relevant to the query by recognizing from the 'Morgan District Brigade' table segment that the 'Disbanded' column value was 1782. 
However, information related to this query was present in two passages: '2nd Rowan County Regiment' and 'Burke County, North Carolina'. 
The LLM incorrectly verified only 'Burke County, North Carolina' as relevant, likely due to its more plausible-sounding title, while overlooking the correct answer 'North-Carolina' in the passage titled '2nd Rowan County Regiment'.
Consequently, the system produced an incorrect response, 'North Carolina'.
This error highlights two problems: (i) a limitation of the LLM reasoning capability and (ii) an example case of the \texttt{OTT-QA} benchmark's wrong answer annotation.


Figure 6(d) shows another failure case of the passage verification module, this time for the query, "When was the first album of Travie McCoy's discography that he guest appeared on?".
Prior retrieval results correctly introduced the ground truth table titled 'Travie McCoy discography (Guest Appearances)' to the passage verification module.
However, the LLM incorrectly inferred that "the table does not specify the information about Travie McCoy" as seen in the second line of its response bubble.
It then relied on its parameterized knowledge to wrongly verify a passage titled 'This Is How It Goes Down as relevant'. 
The correct answer 'Funhouse (Pink album)' was excluded from the final retrieved document set due to the verification error.

\section{Software and Data Licenses}
The licenses for the software and datasets used in this paper are as follows:
\begin{itemize}
    \item ColBERTv2: MIT License
    \item bge-reranker-v2-minicpm-layerwise: Apache 2.0 License
    \item LLaMA 3.1-8B-Instruct: LLaMA 3.1
    \item LLaMA 3.1-70B-Instruct: LLaMA 3.1
    \item Chain-of-Table: Apache 2.0 License
    \item TableLlama: MIT License
    \item COS: GPL-3.0 License
    \item OTT-QA: MIT License
\end{itemize}
All software and datasets were used strictly for research purposes and were not utilized in any non-research contexts, particularly for commercial applications.
The datasets used in this study, OTT-QA and MultimodalQA, consist of English-language data sourced from the Wikipedia domain.

\section{AI Assistants}
We used ChatGPT-4o (\citealp{hurst2024gpt}) to debug code efficiently, quickly identifying and resolving errors in our implementations. 
Additionally, we used it for rephrasing sentences in my writing to improve clarity and readability.




\section{Reproducibility Statement} 
OTTeR (\citealp{OTTeR}) and DoTTeR (\citealp{DoTTeR}) were reproduced using the official code available at \href{https://github.com/Jun-jie-Huang/OTTeR}{\texttt{OTTeR}} and \href{https://github.com/deokhk/DoTTeR}{\texttt{DoTTeR}}, respectively. COS (\citealp{COS}) and CORE (\citealp{CORE}) were reproduced using the official code from \href{https://github.com/Mayer123/UDT-QA}{\texttt{UDT-QA}}. 
The source code, data, and other artifacts for \texttt{HELIOS} have been made available at \href{https://github.com/pshlego/HELIOS}{\texttt{HELIOS}}.

\section{Prompt Templates}
\label{sec:Prompts}

For Star-based LLM refinement, we extended the prompt from \texttt{Chain-of-Table} (\citealp{chainoftable}), originally designed for selecting relevant rows from tables, to support column-wise aggregation and passage verification. 
This extension enables the joint consideration of table segments and linked passages. 
For the LLM reader, we constructed the prompt based on the HybridQA prompt from \texttt{TableLlama} (\citealp{tablellama}).















\newpage
\newpage

% \subsection{Prompt for Aggregation Query Classification}
% \label{sec:aggregation_query_detection_prompt}

\begin{figure*}[htb]

\begin{tcolorbox}
[title = Aggregation Query Classification, colback = gray!10, colframe = black, sharp corners, boxrule=0.5mm]

Using \texttt{f\_agg()} API, return True to detect when a natural language query involves performing aggregation operations (max, min, avg, group by). 
Strictly follow the format of the below examples. Please provide your explanation first, then answer the question in a short phrase starting by ’Therefore, the answer is:’ \\
 \\
\textbf{Question}: when was the third highest paid Rangers F.C. player born? \\
\textbf{Explanation}: The question involves finding the birth date of the third highest paid player, which requires aggregation to find the third highest paid player. Therefore, the answer is: \texttt{f\_agg([True])} \\

\textbf{Question}: what is the full name of the Jesus College alumni who graduated in 1960? \\
\textbf{Explanation}: The question involves finding the full name of the alumni who graduated in 1960, which does not require aggregation. Therefore, the answer is: \texttt{f\_agg([False])} \\

\textbf{Question}: how tall, in feet, is the Basketball personality that was chosen as MVP most recently? \\
\textbf{Explanation}: The question involves finding the most recent MVP winner, which requires aggregation to identify the relevant player. Therefore, the answer is: \texttt{f\_agg([True])} \\

\textbf{Question}: what is the highest best score series 7 of Ballando con le Stelle for the best dancer born 3 July 1969? \\
\textbf{Explanation}: The question involves finding the highest score in a series for a specific dancer, which requires aggregation. Therefore, the answer is: \texttt{f\_agg([True])} \\

\textbf{Question}: which conquerors established the historical site in England that attracted 2,389,548 2009 tourists? \\
\textbf{Explanation}: The question involves identifying the conquerors who established a historical site, which does not require aggregation. Therefore, the answer is: \texttt{f\_agg([False])} \\

\textbf{Question}: what is the NYPD Blue character of the actor who was born on January 29, 1962? \\
\textbf{Explanation}: The question involves finding the character played by an actor born on a specific date, which does not require aggregation. Therefore, the answer is: \texttt{f\_agg([False])} \\
 \\
\textbf{Question}: `{\texttt{\{question\}}}' \\
\textbf{Explanation}: 

\end{tcolorbox}
\label{fig:agg_detection}
\end{figure*}









\newpage
\newpage

% \subsection{Prompt for Column-wise Aggregation}
% \label{sec:table_wise_aggregation_prompt}

\begin{figure*}[htb]

\begin{tcolorbox}
[title = Column-wise Aggregation, colback = gray!10, colframe = black, sharp corners, boxrule=0.5mm]

Using \texttt{f\_row()} API to select relevant rows in the given table and linked passages that support or oppose the question. 
Strictly follow the format of the below example. Please provide your explanation first, then select relevant rows in a short phrase starting by: \textit{``Therefore, the relevant rows are:"}
\\
\\
\texttt{/* 
table caption : list of rangers f.c. records and statistics \\
col : \# | player | to | fee | date \\
row 1 : 1 | alan hutton | tottenham hotspur | 9,000,000 | 30 january 2008 \\
row 2 : 2 | giovanni van bronckhorst | arsenal | 8,500,000 | 20 june 2001 \\
row 3 : 3 | jean-alain boumsong | newcastle united | 8,000,000 | 1 january 2005 \\
row 4 : 4 | carlos cuellar | aston villa | 7,800,000 | 12 august 2008 \\
row 5 : 5 | barry ferguson | blackburn rovers | 7,500,000 | 29 august 2003 
*/}
 \\
\texttt{/* 
Passages linked to row 1: \\
- Alan Hutton: Alan Hutton (born 30 November 1984) is a Scottish former professional footballer, who played as a right back. Hutton started his career with Rangers, and won the league title in 2005. \\
- Tottenham Hotspur F.C.: Tottenham Hotspur Football Club, commonly referred to as Tottenham or Spurs, is an English professional football club in Tottenham, London, that competes in the Premier League.
*/}
 \\
\texttt{/* 
Passages linked to row 2: \\
- Giovanni van Bronckhorst: Giovanni Christiaan van Bronckhorst (born 5 February 1975), also known by his nickname Gio, is a retired Dutch footballer and currently the manager of Guangzhou R\&F. 
*/}
 \\
\texttt{/* 
Passages linked to row 3: \\
- Jean-Alain Boumsong: Jean-Alain Boumsong Somkong (born 14 December 1979) is a former professional football defender, including French international.  \\
- Newcastle United F.C.: Newcastle United Football Club is an English professional football club based in Newcastle upon Tyne, Tyne and Wear, that plays in the Premier League, the top tier of English football.
*/}
 \\
\textbf{Question:} '\texttt{When was the third highest paid Rangers F.C . player born ?}' \\
\textbf{Explanation:} The third-highest paid Rangers F.C. player, Jean-Alain Boumsong (row 3). \textit{Therefore, the relevant rows are:} \texttt{f\_row([row 3])}'
 \\
 \\
\texttt{/* '{\texttt{\{table\}}}' */} \\
 \\
\texttt{/* '{\texttt{\{linked\_passages\}}}' */} \\
 \\
\textbf{Question:} '{\texttt{\{question\}}}' \\
\textbf{Explanation:}

\end{tcolorbox}
\label{fig:column_aggregation}
\end{figure*}













\newpage
\newpage

% \subsection{Prompt for Passage Verification}
% \label{sec:passage_wise_selection_prompt}

\begin{figure*}[htb]

\begin{tcolorbox}
[title = Passage Verification, colback = gray!10, colframe = black, sharp corners, boxrule=0.5mm]


Using \texttt{f\_passage()} API to return a list of passage titles that are relevant to the question, even if they are only partially related. 
Strictly follow the format of the below example. Please provide your explanation first, then return a list of passages in a short phrase starting by: \textit{``Therefore, relevant passages are:"}
 \\
  \\
\texttt{/* 
table caption : List of politicians, lawyers, and civil servants educated at Jesus College, Oxford \\
col : Name | M | G | Degree | Notes \\
row 1 : Lalith Athulathmudali | 1955 | 1960 | BA Jurisprudence (2nd, 1958), BCL (2nd, 1960) | President of the Oxford Union (1958); a Sri Lankan politician; killed by the Tamil Tigers in 1993 
*/}
 \\
\texttt{/* 
List of linked passages: ["Law degree", "Oxford Union", "Lalith Athulathmudali"] \\
Title: Lalith Athulathmudali. Content: Lalith William Samarasekera Athulathmudali, PC (Sinhala; 26 November 1936 - 23 April 1993), known as Lalith Athulathmudali, was a Sri Lankan statesman. He was a prominent member of the United National Party, who served as Minister of Trade and Shipping; Minister of National Security and Deputy Minister of Defence; Minister of Agriculture, Food and Cooperatives, and finally Minister of Education. \\
Title: Law degree. Content: A law degree is an academic degree conferred for studies in law. Such degrees are generally preparation for legal careers; but while their curricula may be reviewed by legal authority, they do not themselves confer a license. A legal license is granted (typically by examination) and exercised locally; while the law degree can have local, international, and world-wide aspects. \\
Title: Oxford Union. Content: The Oxford Union Society, commonly referred to simply as the Oxford Union, is a debating society in the city of Oxford, England, whose membership is drawn primarily from the University of Oxford. Founded in 1823, it is one of Britain’s oldest university unions and one of the world’s most prestigious private students' societies. The Oxford Union exists independently from the university and is separate from the Oxford University Student Union.
*/}
 \\
  \\
\textbf{Question:} What is the full name of the Jesus College alumni who graduated in 1960? \\
\textbf{Explanation:} First, Lalith Athulathmudali graduated in 1960. Second, the linked passage titled ``Lalith Athulathmudali" confirms his full name. \textit{Therefore, relevant passages are:} \texttt{f\_passage(["Lalith Athulathmudali"])}
 \\
 \\
\texttt{/* '{\texttt{\{table\}}}' */} \\
 \\
\texttt{/* `{\texttt{\{linked\_passages\}}}' */} \\
 \\
\textbf{Question:} `{\texttt{\{question\}}}' \\
\textbf{Explanation:}

\end{tcolorbox}

\end{figure*}


\newpage
\begin{figure*}[htb]

\begin{tcolorbox}
[title = Hybrid Question Answering, colback = gray!10, colframe = black, sharp corners, boxrule=0.5mm]
Below is an instruction that describes a task, paired with an input that provides further context. Write a response that appropriately completes the request.
\\
\#\#\# Instruction: \\
This is a hybrid question answering task. \\
The goal of this task is to answer the question given tables and passages. \\
Strictly follow the format of the below examples. Please provide your explanation first, then answer the question in a short phrase starting by: \textit{`Therefore, the answer is:`}
 \\
\#\#\# Examples: \\
\texttt{
Title : List of politicians, lawyers, and civil servants educated at Jesus College, Oxford \\
col : Name | M | G | Degree | Notes \\
row 1 : Lalith Athulathmudali | 1955 | 1960 | BA Jurisprudence ( 2nd , 1958 ) , BCL ( 2nd , 1960 ) | President of the Oxford Union ( 1958 ) ; a Sri Lankan politician ; killed by the Tamil Tigers in 1993 \\
row 2 : Neal Blewett ( HF ) | 1957 | 1959 | BA PPE ( 2nd ) | Member of the Australian House of Representatives ( 1977-1994 ) , Government Minister ( 1983-1994 ) , High Commissioner to the UK ( 1994-1998 ) \\
row 3 : Joseph Clearihue | 1911 | 1914 | BA Jurisprudence ( 2nd , 1913 ) , BCL ( 3rd , 1914 ) | Canadian Rhodes scholar ; later became a member of the Legislative Assembly of British Columbia and a county court judge ; also chairman of the council of Victoria College , British Columbia ( which became the University of Victoria under his leadership ) \\
}
 \\
\texttt{ 
Passages linked to row 1: \\
- [Lalith Athulathmudali](https://en.wikipedia.org/wiki/Lalith\_Athulathmudali) Lalith William Samarasekera Athulathmudali , PC (26 November 1936 - 23 April 1993) , known as Lalith Athulathmudali , was Sri Lankan statesman . He was a prominent member of the United National Party , who served as Minister of Trade and Shipping ; Minister National Security and Deputy Minister of Defence. \\
Passages linked to row 3: \\
- [Joseph Clearihue](https://en.wikipedia.org/wiki/Joseph\_Clearihue) Joseph Badenoch Clearihue ( 1887-1976 ) was a Canadian lawyer , judge , academic and politician . \\
}
\\
\textbf{Question:}What is the full name of the Jesus College alumni who graduated in 1960? \\
\textbf{Explanation:} Lalith Athulathmudali graduated in 1960, and his full name is Lalith William Samarasekera Athulathmudali. Therefore, the answer is: Lalith William Samarasekera Athulathmudali.
\\
\#\#\# Input: \\
Here are the table and passages to answer this question. Please provide your explanation first, then
answer the question in a short phrase starting by \textit{`Therefore, the answer is:`}
\\
 \\
\texttt{/* `{\texttt{\{table\_and\_linked\_passages\}}}' */} \\
 \\
\textbf{Question:} `{\texttt{\{question\}}}' \\
\textbf{Explanation:}

\end{tcolorbox}
\label{fig:passage_verification}
\end{figure*}